\documentclass{article}
\PassOptionsToPackage{numbers}{natbib}
\usepackage[preprint]{neurips_2026}
\usepackage[utf8]{inputenc}
\usepackage[T1]{fontenc}
\usepackage{hyperref,url,booktabs,amsmath,amssymb,microtype}
\title{Predictive evaluation with limited labels:\\
independent replication, audit, and conditional cross-subject shrinkage}
\author{Francisco Angulo de Lafuente\\Independent researcher, Madrid, Spain\\
\texttt{lareliquia.angulo@gmail.com}}
\begin{document}
\maketitle

\begin{abstract}
We independently audit and extend a lightweight submission for the Predictive
AI Evaluation Competition. Reconstructing the original sampling design from
the released response tables reproduces its empirical-mean reference score
of 0.2453845 and Beta--Binomial score of 0.2104970 to numerical precision.
The strongest saved alternatives were already disclosed; we find no observable
evidence that their better point estimates were concealed. We do find problems
in the isolation of development data for richer models, eligibility summaries,
and the scope of previous negative conclusions. A new evaluator uses all
13,200 eligible subject--benchmark pairs from 151 public benchmarks, eight
acquisition/evaluation splits, grouped nested validation, and component-level
uncertainty. It confirms an advantage for the original shrinkage rule:
0.236814 to 0.205984 against the empirical mean. Three new isolated-evidence
families do not establish a reliable improvement. A separately registered
extension uses only supplied labels of other subjects from the same benchmark.
With simulated four-subject cohorts, a fallback-preserving variant improves
0.206027 to 0.203465, with a 99\% component-bootstrap interval for the
difference of $[-0.006257,-0.000551]$. With no peer evidence it is exactly
the original predictor. This result is conditional on evidence availability;
neither cohort exposure nor private-test performance has been verified on
Codabench. All candidates and failed approaches are retained.
\end{abstract}

\section{Task, estimand, and scope}
The competition asks for the probability of an observed binary outcome for an
AI subject and item, given visible attributes and a limited set of revealed
labels~\cite{competition}. Subjects are configurations, including the model,
harness, and reasoning settings. Local benchmark names and internal subject
IDs are used to construct evaluation partitions; the submitted predictor uses
only the supplied dictionaries and anonymous benchmark identifiers.

At budgets $n\in\{0,1,3,7,15,31\}$, let $B_n$ be the mean squared error over
all recorded evaluation responses of a subject--benchmark pair. Its score is
\begin{equation}
 L=0.1B_0+0.2B_1+0.2B_3+0.2B_7+0.2B_{15}+0.1B_{31}.
\end{equation}
The final estimand averages $L$ equally across pairs, irrespective of their
number of responses. Repeated observations are retained within each pair;
they do not create additional equally weighted pairs. Lower scores are better.
Acquisition and evaluation items are disjoint and budgets are nested.

Our results are local revalidation on a public collection. The public
benchmarks had been inspected during an earlier sprint; the new partitions
are not previously unobserved scientific data. A new evaluation program and
a reviewer who did not inspect the new scores provide checks against shared
implementation errors, not a guarantee against all forms of adaptive analysis.
No local score in this paper is an official private-test score or ranking.

\section{Independent reconstruction and provenance}
The audit read the fifteen original analysis scripts, fourteen released
result files, report, model, and retained development logs. Its principal
evaluator imports none of those analysis modules and does not deserialize
the original \texttt{pares.pkl}. It reads the original Parquet tables in
bounded batches, aggregates counts, successes, and the first observation
for each subject--item combination, and independently constructs eligibility
and acquisition trajectories.

\paragraph{Frozen inputs.}
The local cache records the following dataset revision:
\begin{center}\small
\texttt{54074453da7b9911be0f48e4c560be547bc84d21}
\end{center}
All 278 public response
files match their Hugging Face SHA-256 download records; 501 supporting item,
subject, and benchmark files were also hashed and match. The release manifest
pins 890 necessary public files, including all public benchmark metadata,
with a combined size of 5,354,371,887 bytes. The public baseline source
matches upstream after newline normalization and has an identical Python AST.
Data are obtained from the permitted measurement database~\cite{data}; no
outside training data or paid inference service was used.

The old report cited an earlier revision. A historical manifest comparison
found 32 public response paths absent from that older snapshot and one changed
response file. This makes the old revision an insufficient reproduction
instruction. It does not by itself invalidate the reported numbers: the
14,733 previously aggregated subject--benchmark rows match the independently
recomputed successes and response counts exactly. Two additional all-zero
pairs contain 116 responses in total and are ineligible because they have
60 and one distinct items. The old binary filter excluded them. The current
binary scan contains 70,193,656 physical rows, of which 68,840,604 have an
observed response; missing values must not be counted as labels.

\paragraph{Exact replication of the earlier design.}
Separately from the primary evaluation, we recreated the old six-pair and
150-evaluation-item sampling rules. Source row order was recovered directly
from the response tables, rather than inferred from a hash-group output
ordering. All 769 sampled pairs were reconstructed, including all 496 pairs
in the original 101-benchmark headline set. Recomputed pair scores differ
from the saved Beta--Binomial and empirical-mean scores by at most
$1.11\times10^{-16}$. Their means are respectively 0.2104969827 and
0.2453845123. Thus the 14.2\% headline gain is a reproducible result for
that particular local design. It does not establish superiority to stronger
competitors or to an unknown private-test distribution.

\section{Audit findings and corrections}
\paragraph{Better saved results were disclosed.}
The earlier report lists a subject-conditioned empirical-Bayes predictor at
0.20989535 and a macro hierarchical predictor at 0.21024007, both below the
submitted predictor's point estimate. The former also heads the README table.
Their uncertainty and runtime simplicity motivated the earlier choice.
The retained artifacts provide no evidence of concealment. An audit cannot
recover deleted or unsaved experiments or establish an author's intent.

\paragraph{Development isolation for richer priors was incomplete.}
The earlier LOBO procedure excluded the target benchmark but fitted a richer
prior on every other benchmark, even during development selection. Its
nominal headline set therefore contributed labels to those fitted priors.
This is a legitimate transductive evaluation only if explicitly described
as such; it is not a fully isolated development/headline partition. The
constant Beta--Binomial rule fits no prior from those aggregates and does
not directly inherit that fitted-prior defect. The new evaluator separates
fitting, inner validation, and external testing by benchmark components.

\paragraph{Eligibility and weighting require distinct items.}
Ninety public pairs have at least 80 recorded responses but fewer than
80 distinct items. Some old descriptive statistics and empirical-prior filters
counted them as eligible. On the old headline benchmarks, response-based
counting gives 11,657 pairs; the correct distinct-item criterion gives
11,639. The full new population contains 13,200 eligible pairs. The old
cap of six subjects per benchmark changes the estimand and cannot be repaired
exactly by merely reweighting those six observations. We remove both that cap
and the 150-item evaluation cap.

\paragraph{Negative text experiments were narrower than their conclusions.}
The program called TF--IDF used word and character hashing without inverse
document frequency. It predicted an item residual, evaluated against an
in-sample empirical pair mean, with limited training rows and one optimizer
configuration. Another text model used response-weighted samples and a
budget-zero endpoint. These experiments are informative failures of those
specific configurations; they do not establish that transferable item signal
is absent, that pretrained representations cannot help, or that future
improvement requires paid runtime language-model calls. We do not relabel
them as independent successful replications in this audit.

\paragraph{Uncertainty and oracle language.}
The old interval $[-0.0373,-0.0249]$ belongs to a randomized acquisition run,
not the headline trajectory. The original report distinguishes them. An
independent 10,000-benchmark bootstrap on the headline artifacts gives
$[-0.04107,-0.02854]$, confirming its advantage under that design.
The older ``sign-flip'' routine was an exact binomial sign test, and zero
opposite-sign bootstrap samples are not a true probability of zero.
Finally, a retrospectively fitted per-item oracle is an empirical lower
score on the observed sample; its within-item variance is not a demonstrated
irreducible noise floor for new responses.

\section{New evaluation and three initial improvements}
For each eligible pair, items are canonically ordered and split 50/50 under
eight SHA-256-derived seeds. One response is drawn uniformly from the recorded
responses of each acquired item; equivalent Bernoulli draws from that item's
empirical success fraction avoid expanding repeated records. All evaluation
responses contribute to scoring. For any predictor constant within a pair,
the exact sufficient-statistic loss is
\begin{equation}
 B(p)=p^2-2pq+q,\qquad q=\frac{\text{evaluation successes}}{\text{evaluation responses}}.
\end{equation}
This identity was checked against explicit squared errors on 150 synthetic
samples. The primary trajectories are invariant to storage row order.

Exact shared content hashes connect the 151 evaluation benchmarks into
142 components. We keep each component within one of five external folds.
Inside each external training set, another deterministic component split
defines fitting and validation roles. Every feature transform and fitted
model sees only the fitting role; every hyperparameter choice sees only
inner validation. Public population imbalance is retained: for example,
\texttt{alignment\_faking} contributes 4,278 pairs and \texttt{sib200}
contributes 3,341. Folds balance component counts, not pair counts.

The simple controls are the empirical mean and
\begin{equation}
 \widehat p=\frac{k+1.5}{n+3},
\end{equation}
the posterior mean of a symmetric Beta prior. We claim no novelty for this
standard shrinkage rule. The first three predeclared improvements are:
(A) a Beta mean/strength grid; (B) a conditional table indexed by budget and
success count, regularized and mixed with the original predictor; and
(C) a ridge prior using visible subject attributes, with categorical and name
token features. All validation candidates are retained. The table regularizer
has weight sixteen simulated rows, approximately two training pairs across
eight seeds, not sixteen independent subjects.

\begin{table}[t]
\centering\small
\caption{Independent isolated-pair evaluation. All 13,200 eligible pairs and
eight seeds; no cross-subject labels. Deltas compare with original Beta--Binomial.
Intervals are component bootstrap intervals, averaging seeds before resampling.}
\label{tab:initial}
\begin{tabular}{lrrl}
\toprule
Method & Brier ALC & Delta & 95\% interval for delta\\
\midrule
Constant 0.5 & 0.250000 & & \\
Empirical-mean reference & 0.236814 & & \\
Original Beta--Binomial & 0.205984 & & \\
A: tuned Beta prior & 0.205535 & $-0.000449$ & $[-0.005350,+0.003663]$\\
B: calibration table & 0.205877 & $-0.000107$ & $[-0.000522,+0.000105]$\\
C: subject ridge prior & 0.216549 & $+0.010565$ & $[+0.001218,+0.017943]$\\
\bottomrule
\end{tabular}
\end{table}

The original shrinkage rule beats the empirical mean by 0.030830 ALC,
with a 95\% component-bootstrap interval of $[-0.038426,-0.027214]$ and an
improvement on 133 of 151 benchmarks. Table~\ref{tab:initial} shows why the
small point improvements of A and B do not justify promoting either. C
performs worse. Variation across acquisition seeds does not replace
uncertainty about transfer to new benchmarks.

\section{Conditional use of cross-subject evidence}
After inspecting the initial experiments, we separately registered an
exploratory extension motivated by the documented interface: supplied
\texttt{labeled} entries can include different subject--benchmark pairs.
The submitted program can therefore use another visible subject's acquired
labels when their anonymous benchmark matches. It must never use an unseen
evaluation label or assume a persistent shared cache.

For a target subject let $(k,n)$ summarize its matching visible attributes,
and let $(k_o,n_o)$ summarize other supplied subjects on the same benchmark.
Our fallback-preserving predictor uses
\begin{equation}
 m_o=\frac{1.5+k_o}{3+n_o},\quad
 m=0.5+\rho(m_o-0.5),\quad
 \widehat p=\frac{k+3m}{n+3}.
\end{equation}
With no peer evidence $m_o=m=0.5$, so this is exactly the original predictor.
A variant that also retunes strength is retained as a separate failed
candidate: it loses this protective property and worsens isolated-pair loss.

We simulate deterministic random cohorts of four subjects per benchmark as
the primary availability scenario, with sizes one, two, and eight as
sensitivities. The final partial block can have fewer subjects. Both methods
receive the same cohort's acquired labels at each budget. Eight visible
attribute signatures are shared by 59 internal subject pairs; matching
dictionary attributes, rather than inaccessible IDs, correctly pools those
labels for both the original reference and the extension. Thus the cohort
reference differs slightly from the isolated-ID simulation in
Table~\ref{tab:initial}; they must not be conflated.

\begin{table}[t]
\centering\small
\caption{Conditional availability experiments. The same configuration selected
in inner validation with cohorts of four is used in the other scenarios.
These are separate simulated environments, not a choice of the best environment.}
\label{tab:context}
\begin{tabular}{rrrrl}
\toprule
Subjects/cohort & Original API & Fallback extension & Delta & 95\% interval\\
\midrule
1 & 0.205984 & 0.205984 & $0$ & exact probability equality\\
2 & 0.206010 & 0.204711 & $-0.001298$ & $[-0.003207,-0.000152]$\\
4 & 0.206027 & 0.203465 & $-0.002562$ & $[-0.005234,-0.000938]$\\
8 & 0.206030 & 0.202565 & $-0.003465$ & $[-0.006599,-0.001555]$\\
\bottomrule
\end{tabular}
\end{table}

For cohorts of four, the relative reduction is 1.24\% against its correctly
matched API reference, favorable in all eight acquisition seeds. Adjusting
for the three initial candidates and the two contextual candidates gives
a 99\% interval of $[-0.006257,-0.000551]$. With retuned strength instead,
ALC is 0.207251, a worse point estimate than the cohort reference. We retain
this negative result. The exported fallback candidate uses $\rho=0.5$,
chosen by the pair-count-weighted mean of logged inner-validation scores
at fixed strength three. The headline contextual result estimates the
nested selection procedure; it is not a fresh unbiased test of a fixed
parameter chosen after fitting to all public development data.

The simulator-to-function correspondence was checked independently on
5,376 synthetic cases across cohort sizes, mixtures, strengths, visible
collisions, and partial blocks: maximum error zero. Another 528 valid
$(n,k)$ states confirm exact fallback. Probability-level equality is tested
exactly; mathematically identical score reductions are canonicalized only
after a $10^{-14}$ numerical check to avoid mistaking floating-point noise
for a nonzero effect.

\section{Inference, limitations, and reproducibility}
We average the eight seeds within each pair, then bootstrap 10,000 complete
content components. Each replicate recomputes total score difference divided
by total pair count, preserving the official pair-weighted estimand rather
than silently averaging benchmark means. Reported 95\% intervals are
descriptive; family-adjusted intervals are used when comparing promoted
candidate families. The grouped sign-randomization calculation uses a
plus-one correction and requires symmetry assumptions; it is not treated
as an exact distribution-free statement about transfer.

This bootstrap conditions on fitted cross-validation models; it does not
retrain the whole selection process for each replicate. Exact content hashes
do not catch paraphrases or every family relationship. Public-population
weights may differ from the platform's subject sampling, and binary outcomes
have heterogeneous meanings across released benchmarks. The inner validation
sets and folds can be imbalanced by a few large benchmarks. Fitted models
are not refitted on fit plus validation after selection, which is valid but
may sacrifice efficiency. We make no claim to exhaust model classes or
the potential value of pretrained representations.

Most importantly, Codabench was inaccessible during this audit. The interface
allows peer labels, but their actual simultaneous availability, ordering,
and cohort sizes are unverified. The contextual result is therefore a
conditional local improvement and a candidate for platform testing. When no
such labels are supplied, its output is exactly the original rule. The
original ZIP is retained alongside the candidate. The development release is
public, but no Codabench submission ID or official score has been obtained.

The accompanying release contains both single-file predictors, the original
artifacts, the independent evaluator, grouped split manifests, all inner
validation candidates, per-pair scores, source hashes, protocols registered
before their experiments, and the independent review. A pinned downloader
replaces the earlier unversioned download helper. Inputs remain subject to
their upstream access conditions and licenses and are not bundled as raw
content. Python package versions and reproduction commands are recorded in
the release README. Tests and experiments used CPU with at most four library
threads; neither GPU nor paid inference was necessary for these results.

Code, protocols, results, and reproduction instructions are publicly available
at \url{https://github.com/Agnuxo1/AIMS-PAIEC-2026}; the presentation is linked
from that repository. The immutable initial release is tagged \texttt{v1.0.0}.
An accepted Codabench submission ID remains pending because the platform is
unreachable. The organisers were contacted on 24 October 2026 about the
submission route; no alternative acceptance has been received. The current
report is a technical preprint, not evidence of admission or a presentation award.

\paragraph{Disclosure.}
Large language models assisted the original development and this audit,
including code, writing, and an independent review conducted without viewing
the new experimental scores. The new data reconstruction and scoring were
executed with deterministic tools, and intermediate failures are retained:
one replication attempt incorrectly assumed hash-group output order, and
two checks required technical fixes for JSON serialization and floating-point
reduction. These were corrected without changing candidate grids or removing
unfavorable results. No raw credential file was inspected or included.

\begin{thebibliography}{3}
\bibitem{competition} Stanford AIMS. Predictive AI Evaluation Competition at NeurIPS 2026.
\url{https://aims.stanford.edu/competition/}. Accessed 24 October 2026.
\bibitem{data} AIMS Foundations. Measurement database, frozen revision
\texttt{54074453}. \url{https://huggingface.co/datasets/aims-foundations/measurement-db}.
\bibitem{baseline} AIMS Foundations. PAIEC predictor baselines.
\url{https://github.com/aims-foundations/paiec_baseline}.
\end{thebibliography}
\end{document}
